\documentclass[12pt]{amsart} 
\usepackage{mathtools}

\input{./latex-preamble/cm-preamble.tex}
\input{./latex-preamble/use-scr.tex}
\input{./latex-preamble/fr-mod.tex}

%\graphicspath{ {./diagrams/} }

\begin{document}

{Hartshorne} \hfill {9-23-2026}

\bigskip\bigskip

\begin{center}

{\sc Chapter III.8}

\end{center}

\begin{center}

  {Noah Parker} 
    
\end{center}

\bigskip\bigskip

\begin{enumerate}
  \item[8.1.] {
      Let $f:X\to Y$ be a continuous map of topological spaces.
      Let $\F$ be a sheaf of abelian groups on $X$, and assume that $R^if_*(\F)=0$ for all $i>0$.
      Show that there is a natural isomorphism
      \[
        H^\bullet(X,\F)\cong H^\bullet(Y,f_*\F).
      \] 
      (This is a degenerate case of the Leray spectral sequence.)
    }
\end{enumerate}

\begin{proof}
  Note that we may regard $X$ as a ringed space by setting $\O_X=\ul \ZZ$, so that $\Ab(X)\cong\Mod(X)$.
  We also have 
  \[
    H^0(X,\F)=\Gamma(X,\F)=\Gamma(Y,f_*\F)= H^0(Y,f_*\F).
  \]
  Thus, it suffices to show that $H^\bullet(Y,f_*(-))$ is a universal $\delta$-functor.
  $f_*$ is exact by assumption, so it suffices by (1.3A) to show that $H^i(Y, f_*(-))$ is effaceable for each $i>0$.
  Each injective object $\I$ of $\Ab(X)$ is flasque by (2.4), so each $f_*\I$ is flasque by (II, Ex. 1.16d).
  Then $H^i(f_*\I)=0$ for all $i>0$ by (2.5), so we're done.
  % Take an injective resolution $0\to \F\to \I^\bullet$ for $\F$.
  % Then each $\I^i$ is flasque by (2.4), so each $f_*\I^\bullet$ is flasque as well by (Ex. II.1.16d).
  % $f_*$ is exact since $R^if_*(\F)$, so $0\to f_*\F\to f_*\I^\bullet$ is a flasque resolution of $\F$.
  % It follows that $H^\bullet(Y,f_*\F)\cong h^\bullet(f_*\I^\bullet)$.
  % Clearly $H^\bullet(X,\F)\cong h^\bullet(\I^\bullet)$, so
\end{proof}

\begin{enumerate}
  \item[8.2.] {
      Let $f:X\to Y$ be an affine morphism of schemes with $X$ noetherian, and let $\F$ be a quasi-coherent sheaf on $X$.
      Show that the hypotheses of (Ex. 8.1) are satisfied, and hence that $H^\bullet(X,\F)\cong H^\bullet(Y,f_*\F)$.
      (This gives another proof of (Ex. 4.1).)
    }
\end{enumerate}

\begin{proof}
  It suffices to show that $R^if_*(\F)=0$ for all $i>0$.
  Take $U\subset Y$ an affine open, so that $V:=f^{-1}(U)$ is affine.
  (8.2) and (8.5) imply
  \[
    R^\bullet f_*(\F)|_U\cong R^\bullet(f|_V)_*(\F|_V)\cong  H^\bullet(V,\F)^\sim,
  \] 
  whence
  \[
    R^i f_*(\F)|_U\cong H^i(V,\F)^\sim = 0
  \] 
  for all $i>0$ by (3.5).
  It follows that $R^if_*(\F)=0$ for all $i>0$, as desired.
\end{proof}

\begin{enumerate}
  \item[8.3.] {
      Let $f:X\to Y$ be a morphism of ringed spaces, let $\F$ be an $\O_X$-module, and let $\E$ be a locally free $\O_Y$-module of finite rank.
      Prove the \emph{projection formula} (cf. (II, Ex. 5.1))
      \[
        R^\bullet f_*(\F\otimes f^*\E)\cong R^\bullet f_*(\F)\otimes \E.
      \] 
    }
\end{enumerate}

\begin{proof}
  For $i=0$, this is precisely the projection formula (II, Ex. 5.1d).
  $f^*$ is right adjoint and thus commutes with inverse limits, and in particular finite direct sums.
  $f^*(\O_X)=\O_Y$, so this gives us $f^*\E$ locally free of finite rank.
  Then both $R^\bullet f_*(-\otimes f^*\E)$ and $R^\bullet f_*(-)\otimes \E$ are $\delta$-functors, so it suffices to show each is effaceable.

  Clearly the latter is effaceable.
  Now, take and embedding $\F\inc \I$ of $\F$ into an injective object of $\Mod(X)$.
  (II, Ex. 5.1abc) gives us an isomorphism
  \begin{align*}
    \Hom(-, \I\otimes f^*\E)
      &\cong \Hom(-\otimes (f^*\E)^\vee, \I),
  \end{align*}
  where the latter functor is exact since $(f^*\E)^\vee$ is locally free of finite rank and $\I$ is injective.
  Then $\I\otimes f^*\E$ is injective, whence $R^\bullet f_*(-\otimes f^*\E)$ is effaceable.
  The result follows from (1.3A).
\end{proof}

\end{document}
